\backhead{FURTHER EXERCISES}{khred}
\begin{exlist}
\item Fresh manure has $k=0.05$ d$^{-1}$ at 20 $^\circ$C. Find its half-life at 20 $^\circ$C and at 10 $^\circ$C ($Q_{10}=2$).\exans{13.9 d; 27.7 d}
\item Clover residues (C:N 15) of 3 t DM\,ha$^{-1}$ (42 per cent C) decompose fully with $e=0.4$ and $CN_m=8$. Find the net N released.\exans{21 kg N\,ha$^{-1}$}
\item A herbicide has $K_{oc}=100$ L\,kg$^{-1}$; the soil has 1.5 per cent organic carbon, $\rho_b=1.4$ kg\,L$^{-1}$, $\theta=0.25$ and drains 250 mm\,yr$^{-1}$. Find $K_d$, $R$ and the time to travel 1 m.\exans{1.5 L\,kg$^{-1}$; 9.4; 9.4 yr}
\item A field drains 180 mm at 15 mg NO$_3$-N\,L$^{-1}$. Find the N leached and the mean nitrate concentration in mg NO$_3$\,L$^{-1}$.\exans{27 kg N\,ha$^{-1}$; 66 mg NO$_3$\,L$^{-1}$}
\item Nitrification follows first-order kinetics with $k=0.1$ d$^{-1}$. How long does it take to nitrify 90 per cent of an ammonium application?\exans{23 d}
\end{exlist}
\backhead{ESSAY QUESTIONS}{khred}
{\small\itshape Answer each question in about one page, with equations, a sketch and at least one number from this chapter.}\par\vspace{3pt}
\begin{exlist}
\item Explain the pools of organic matter in Daisy and the meaning of their decay rates.
\item Discuss mineralisation and immobilisation and the role of the C:N ratio of residues.
\item Explain the processes of nitrate loss (leaching, denitrification) and how irrigation affects each.
\item Describe the transport of a pesticide through the soil and the factors that decide whether it reaches the groundwater.
\item Discuss nitrous oxide emission from irrigated soils and how it could be reduced.
\item Explain how the nitrogen balance of a field is closed and checked.
\end{exlist}
